% =================================================
% 文件: 存储.tex
% 章节: 存储技术
% 版本: v1.0
% =================================================

\documentclass[12pt,UTF8]{ctexbook}

\input{../common/page}

\xeCJKsetup{AutoFallBack=true}
\setCJKfamilyfont{hei}{SimHei}
\setCJKfamilyfont{kai}{KaiTi}

\usepackage{titletoc}
\titlecontents{chapter}[0pt]{\vspace{3mm}\bf\addvspace{2pt}\filright}{\contentspush{\thecontentslabel\hspace{0.8em}}}{}{\titlerule*[8pt]{.}\contentspage}

\usepackage[colorlinks,linkcolor=blue,citecolor=blue,urlcolor=blue]{hyperref}

\ctexset{
	part/name= {第,卷},
	part/number={\chinese{part}},
	chapter/name={第,章},
	chapter/number={\arabic{chapter}}
}

\usepackage{graphicx}
\graphicspath{{Images/}}

\usepackage[perpage]{footmisc}

\usepackage{enumitem}

\title{\heiti\zihao{0} 存储技术入门指南}
\author{WangFei}
\date{\today}

\begin{document}

\maketitle
\tableofcontents

\frontmatter
\chapter{前言}

本指南面向初学者，详细介绍存储技术的基础概念、存储类型、文件系统、存储网络、备份恢复等知识。

\mainmatter

\chapter{存储基础概念}
\label{chap:storage_basics}

\section{存储层次}
\label{sec:storage_hierarchy}

存储系统按照速度和容量分为多个层次：

\begin{table}[htbp]
    \centering
    \caption{存储层次}
    \begin{tabular}{|c|c|c|c|}
        \hline
        \textbf{层次} & \textbf{类型} & \textbf{速度} & \textbf{容量} \\
        \hline
        L1 & 寄存器 & 最快 & 最小 \\
        \hline
        L2 & 缓存 & 快 & 小 \\
        \hline
        L3 & 内存 & 较快 & 中等 \\
        \hline
        L4 & SSD & 快 & 较大 \\
        \hline
        L5 & HDD & 慢 & 大 \\
        \hline
        L6 & 磁带 & 最慢 & 最大 \\
        \hline
    \end{tabular}
\end{table}

\section{存储类型}
\label{sec:storage_types}

常见的存储类型：
\begin{table}[htbp]
    \centering
    \caption{存储类型}
    \begin{tabular}{|c|p{8cm}|}
        \hline
        \textbf{类型} & \textbf{说明} \\
        \hline
        直接附加存储 (DAS) & 直接连接到服务器的存储 \\
        \hline
        网络附加存储 (NAS) & 通过网络访问的存储 \\
        \hline
        存储区域网络 (SAN) & 专用网络连接的存储 \\
        \hline
        对象存储 & 通过 HTTP API 访问的存储 \\
        \hline
    \end{tabular}
\end{table}

\section{存储协议}
\label{sec:storage_protocols}

常用的存储协议：
\begin{table}[htbp]
    \centering
    \caption{存储协议}
    \begin{tabular}{|c|p{8cm}|}
        \hline
        \textbf{协议} & \textbf{说明} \\
        \hline
        SCSI & 小型计算机系统接口 \\
        \hline
        iSCSI & IP 网络上的 SCSI \\
        \hline
        Fibre Channel & 光纤通道协议 \\
        \hline
        NFS & 网络文件系统 \\
        \hline
        SMB/CIFS & Windows 文件共享协议 \\
        \hline
        HTTP/REST & 对象存储协议 \\
        \hline
    \end{tabular}
\end{table}

\chapter{硬盘技术}
\label{chap:hard_disk}

\section{HDD}
\label{sec:hdd}

硬盘驱动器（HDD）使用磁性存储：

HDD 的组成：
\begin{itemize}
    \item \textbf{盘片}：存储数据的磁性圆盘
    \item \textbf{磁头}：读写数据的部件
    \item \textbf{主轴}：带动盘片旋转
    \item \textbf{臂}：移动磁头
\end{itemize}

HDD 的性能指标：
\begin{itemize}
    \item \textbf{转速}：每分钟转数（RPM）
    \item \textbf{容量}：存储容量
    \item \textbf{缓存}：内置缓存大小
    \item \textbf{接口}：SATA、SAS
\end{itemize}

\section{SSD}
\label{sec:ssd}

固态硬盘（SSD）使用闪存存储：

SSD 的优点：
\begin{itemize}
    \item 速度快
    \item 无噪音
    \item 低功耗
    \item 抗震性好
\end{itemize}

SSD 的类型：
\begin{table}[htbp]
    \centering
    \caption{SSD 类型}
    \begin{tabular}{|c|p{8cm}|}
        \hline
        \textbf{类型} & \textbf{说明} \\
        \hline
        SATA SSD & SATA 接口 \\
        \hline
        NVMe SSD & PCIe 接口，速度更快 \\
        \hline
        mSATA & 小型 SATA 接口 \\
        \hline
        M.2 & 新型接口，支持 NVMe \\
        \hline
    \end{tabular}
\end{table}

\chapter{RAID}
\label{chap:raid}

\section{RAID 概念}
\label{sec:raid_concept}

RAID（Redundant Array of Independent Disks）是将多个磁盘组合成逻辑卷的技术。

RAID 的级别：
\begin{table}[htbp]
    \centering
    \caption{RAID 级别}
    \begin{tabular}{|c|c|p{4cm}|}
        \hline
        \textbf{级别} & \textbf{特点} & \textbf{用途} \\
        \hline
        RAID 0 & 条带化，无冗余，速度快 & 性能要求高 \\
        \hline
        RAID 1 & 镜像，100\% 冗余 & 数据安全要求高 \\
        \hline
        RAID 5 & 分布式奇偶校验 & 平衡性能和冗余 \\
        \hline
        RAID 6 & 双重奇偶校验 & 更高可靠性 \\
        \hline
        RAID 10 & 镜像+条带 & 高性能+高可靠 \\
        \hline
    \end{tabular}
\end{table}

\section{软件 RAID 配置}
\label{sec:software_raid}

在 Linux 上配置软件 RAID：

创建 RAID 设备：
\begin{verbatim}
mdadm --create /dev/md0 --level=5 --raid-devices=3 /dev/sdb /dev/sdc /dev/sdd
\end{verbatim}

查看 RAID 状态：
\begin{verbatim}
cat /proc/mdstat
mdadm --detail /dev/md0
\end{verbatim}

添加热备盘：
\begin{verbatim}
mdadm --manage /dev/md0 --add-spare /dev/sde
\end{verbatim}

\chapter{LVM}
\label{chap:lvm}

\section{LVM 概念}
\label{sec:lvm_concept}

LVM（Logical Volume Manager）是 Linux 的逻辑卷管理工具。

LVM 的组成：
\begin{itemize}
    \item \textbf{物理卷 (PV)}：物理磁盘或分区
    \item \textbf{卷组 (VG)}：一个或多个物理卷组成
    \item \textbf{逻辑卷 (LV)}：从卷组分配的逻辑分区
\end{itemize}

\section{LVM 配置}
\label{sec:lvm_config}

创建 LVM：
\begin{verbatim}
pvcreate /dev/sdb1 /dev/sdc1
vgcreate vgdata /dev/sdb1 /dev/sdc1
lvcreate -L 50G -n lvdata vgdata
mkfs.xfs /dev/vgdata/lvdata
mount /dev/vgdata/lvdata /data
\end{verbatim}

扩展逻辑卷：
\begin{verbatim}
lvextend -L +20G /dev/vgdata/lvdata
xfs_growfs /data
\end{verbatim}

\chapter{文件系统}
\label{chap:file_systems}

\section{Linux 文件系统}
\label{sec:linux_file_systems}

常见的 Linux 文件系统：
\begin{table}[htbp]
    \centering
    \caption{Linux 文件系统}
    \begin{tabular}{|c|p{8cm}|}
        \hline
        \textbf{文件系统} & \textbf{说明} \\
        \hline
        ext4 & 最常用的文件系统 \\
        \hline
        XFS & 高性能文件系统 \\
        \hline
        Btrfs & 现代文件系统，支持快照 \\
        \hline
        ZFS & 高级文件系统，集成 RAID \\
        \hline
    \end{tabular}
\end{table}

创建文件系统：
\begin{verbatim}
mkfs.xfs /dev/sdb1
mount /dev/sdb1 /data
\end{verbatim}

配置自动挂载：
\begin{verbatim}
echo '/dev/sdb1 /data xfs defaults 0 0' >> /etc/fstab
\end{verbatim}

\section{ZFS 文件系统}
\label{sec:zfs}

ZFS 的特点：
\begin{itemize}
    \item 集成 RAID 功能
    \item 支持快照和克隆
    \item 数据校验
    \item 动态扩展
\end{itemize}

创建 ZFS 池：
\begin{verbatim}
zpool create tank mirror /dev/sdb /dev/sdc
zfs create tank/data
\end{verbatim}

创建快照：
\begin{verbatim}
zfs snapshot tank/data@backup
zfs rollback tank/data@backup
\end{verbatim}

\chapter{存储网络}
\label{chap:storage_network}

\section{NAS}
\label{sec:nas}

网络附加存储（NAS）：

NAS 的特点：
\begin{itemize}
    \item 通过网络访问
    \item 支持文件共享协议
    \item 易于管理
    \item 适合中小企业
\end{itemize}

配置 NFS 服务器：
\begin{verbatim}
dnf install -y nfs-utils
mkdir /data
echo '/data 192.168.1.0/24(rw,sync,no_root_squash)' >> /etc/exports
systemctl enable --now nfs-server
\end{verbatim}

挂载 NFS：
\begin{verbatim}
mount 192.168.1.100:/data /mnt
\end{verbatim}

\section{SAN}
\label{sec:san}

存储区域网络（SAN）：

SAN 的特点：
\begin{itemize}
    \item 专用网络
    \item 块级访问
    \item 高性能
    \item 适合大型企业
\end{itemize}

配置 iSCSI 服务器：
\begin{verbatim}
dnf install -y targetcli
targetcli
/> backstores/block create name=disk1 dev=/dev/sdb
/> iscsi/ create iqn.2024-01.com.example:storage
/> iscsi/iqn.2024-01.com.example:storage/tpg1/luns/ create /backstores/block/disk1
/> iscsi/iqn.2024-01.com.example:storage/tpg1/acls/ create iqn.2024-01.com.example:client
/> saveconfig
\end{verbatim}

配置 iSCSI 客户端：
\begin{verbatim}
dnf install -y iscsi-initiator-utils
iscsiadm -m discovery -t st -p 192.168.1.100
iscsiadm -m node -l
\end{verbatim}

\chapter{备份与恢复}
\label{chap:backup}

\section{备份策略}
\label{sec:backup_strategy}

备份策略类型：
\begin{table}[htbp]
    \centering
    \caption{备份策略}
    \begin{tabular}{|c|p{8cm}|}
        \hline
        \textbf{类型} & \textbf{说明} \\
        \hline
        完全备份 & 备份所有数据 \\
        \hline
        增量备份 & 备份上次备份后变化的数据 \\
        \hline
        差异备份 & 备份上次完全备份后变化的数据 \\
        \hline
    \end{tabular}
\end{table}

3-2-1 备份原则：
\begin{itemize}
    \item 3 份备份
    \item 2 种存储介质
    \item 1 份异地备份
\end{itemize}

\section{备份工具}
\label{sec:backup_tools}

常用备份工具：
\begin{table}[htbp]
    \centering
    \caption{备份工具}
    \begin{tabular}{|c|p{8cm}|}
        \hline
        \textbf{工具} & \textbf{说明} \\
        \hline
        rsync & 文件同步工具 \\
        \hline
        tar & 归档工具 \\
        \hline
        dd & 磁盘镜像工具 \\
        \hline
        Bacula/Bareos & 企业级备份系统 \\
        \hline
        Veeam & 虚拟化备份 \\
        \hline
    \end{tabular}
\end{table}

使用 rsync 备份：
\begin{verbatim}
rsync -av --delete /data/ /backup/data/
\end{verbatim}

使用 tar 归档：
\begin{verbatim}
tar -czf backup.tar.gz /data
\end{verbatim}

\section{数据恢复}
\label{sec:data_recovery}

数据恢复步骤：
\begin{enumerate}
    \item 确认数据丢失范围
    \item 停止使用存储设备
    \item 使用恢复工具
    \item 验证恢复数据
\end{enumerate}

常用恢复工具：
\begin{table}[htbp]
    \centering
    \caption{数据恢复工具}
    \begin{tabular}{|c|p{8cm}|}
        \hline
        \textbf{工具} & \textbf{说明} \\
        \hline
        testdisk & 分区表恢复 \\
        \hline
        photorec & 文件恢复 \\
        \hline
        ddrescue & 磁盘恢复 \\
        \hline
        extundelete & ext 文件系统恢复 \\
        \hline
    \end{tabular}
\end{table}

\chapter{存储性能优化}
\label{chap:performance}

\section{性能指标}
\label{sec:performance_metrics}

存储性能指标：
\begin{itemize}
    \item \textbf{IOPS}：每秒输入输出操作数
    \item \textbf{吞吐量}：每秒数据传输量
    \item \textbf{延迟}：完成一次操作的时间
\end{itemize}

使用 fio 测试性能：
\begin{verbatim}
fio --name=test --size=10G --ioengine=libaio --rw=randread --bs=4k --numjobs=4
\end{verbatim}

\section{优化策略}
\label{sec:optimization_strategy}

存储性能优化：
\begin{itemize}
    \item 使用 SSD 替代 HDD
    \item 合理配置 RAID
    \item 使用缓存
    \item 优化文件系统参数
    \item 使用并行文件系统
\end{itemize}

\chapter{存储监控}
\label{chap:monitoring}

\section{监控工具}
\label{sec:monitoring_tools}

存储监控工具：
\begin{table}[htbp]
    \centering
    \caption{监控工具}
    \begin{tabular}{|c|p{8cm}|}
        \hline
        \textbf{工具} & \textbf{说明} \\
        \hline
        iostat & 磁盘 I/O 统计 \\
        \hline
        vmstat & 虚拟内存统计 \\
        \hline
        smartctl & 磁盘健康状态 \\
        \hline
        Prometheus + Grafana & 综合监控 \\
        \hline
    \end{tabular}
\end{table}

使用 iostat：
\begin{verbatim}
iostat -x 1 10
\end{verbatim}

使用 smartctl：
\begin{verbatim}
smartctl -a /dev/sda
\end{verbatim}

\backmatter

\end{document}